% GENERATED FILE. Edit canonical lessons, then run npm run generate:books.
% canonical-source-hash: fnv1a64:7cc495dbfa786b53
% canonical-lessons: MR-C168-avarne, MR-C168-laksat-yene, MR-C168-lajne, MR-C168-harne, MR-C168-lapne

\chapter{To tidy up, To notice, To be shy, To be defeated, To take cover}
\label{ch:mr-to-tidy-up-to-notice-to-be-shy-to-be-defeated-to-take-cover}
\hlchaptermodality{168}

\begin{chapteropening}
\textbf{By the end of this chapter:} \emph{I can say to tidy up, to notice, to be shy, to be defeated and to take cover.}

Complete the last lesson of chapter 168: I can say to tidy up, to notice, to be shy, to be defeated and to take cover.
\end{chapteropening}

\section[\texorpdfstring{\mr{आवरणे} (āvarṇe)}{आवरणे (āvarṇe)}]{\mr{आवरणे} (āvarṇe) — to tidy up}

\label{lesson:MR-C168-avarne}



\begin{quote}
Before the new one: say the Marathi for to slip, then the Marathi for to suspect.
\end{quote}

\subsection*{You'll want to know: \mr{आवरणे}}
\textbf{\mr{आवरणे}} — \emph{āvarṇe} — \textquotedblleft{}to tidy up\textquotedblright{}.

\begin{grammarlens}[title={What you've built}]
One more word for this chapter.
\end{grammarlens}

\subsection*{Your turn}
\begin{itemize}
\raggedright
  \item \emph{Say it:} \emph{āvarṇe}
  \item \emph{Say it:} \emph{āvarṇe}, once more
  \item \emph{Say it:} say \emph{āvarṇe}
  \item \emph{Recall:} say the Marathi for to slip, then the Marathi for to suspect, then say \emph{āvarṇe} again
\end{itemize}

\begin{tcolorbox}[breakable,colback=teal!4,colframe=teal!35!black,title={Before you move on}]
What does \mr{आवरणे} mean? (\textquotedblleft{}To tidy up\textquotedblright{}.) Say it once more.
\end{tcolorbox}
\section[\texorpdfstring{\mr{लक्षात} \mr{येणे} (lakṣāt yeṇe)}{लक्षात येणे (lakṣāt yeṇe)}]{\mr{लक्षात} \mr{येणे} (lakṣāt yeṇe) — to notice, to realise}

\label{lesson:MR-C168-laksat-yene}



\begin{quote}
Before the new one: say the Marathi for to suspect, then the Marathi for to tidy up.
\end{quote}

\subsection*{You'll want to know: \mr{लक्षात} \mr{येणे}}
\textbf{\mr{लक्षात} \mr{येणे}} — \emph{lakṣāt yeṇe} — \textquotedblleft{}to notice, to realise\textquotedblright{}.

\begin{grammarlens}[title={What you've built}]
One more word for this chapter.
\end{grammarlens}

\subsection*{Your turn}
\begin{itemize}
\raggedright
  \item \emph{Say it:} \emph{lakṣāt yeṇe}
  \item \emph{Say it:} \emph{lakṣāt yeṇe}, once more
  \item \emph{Say it:} say \emph{lakṣāt yeṇe}
  \item \emph{Recall:} say the Marathi for to suspect, then the Marathi for to tidy up, then say \emph{lakṣāt yeṇe} again
\end{itemize}

\begin{tcolorbox}[breakable,colback=teal!4,colframe=teal!35!black,title={Before you move on}]
What does \mr{लक्षात} \mr{येणे} mean? (\textquotedblleft{}To notice, to realise\textquotedblright{}.) Say it once more.
\end{tcolorbox}
\section[\texorpdfstring{\mr{लाजणे} (lājṇe)}{लाजणे (lājṇe)}]{\mr{लाजणे} (lājṇe) — to be shy, to blush}

\label{lesson:MR-C168-lajne}



\begin{quote}
Before the new one: say the Marathi for to tidy up, then the Marathi for to notice.
\end{quote}

\subsection*{You'll want to know: \mr{लाजणे}}
\textbf{\mr{लाजणे}} — \emph{lājṇe} — \textquotedblleft{}to be shy, to blush\textquotedblright{}.

\begin{grammarlens}[title={What you've built}]
One more word for this chapter.
\end{grammarlens}

\subsection*{Your turn}
\begin{itemize}
\raggedright
  \item \emph{Say it:} \emph{lājṇe}
  \item \emph{Say it:} \emph{lājṇe}, once more
  \item \emph{Say it:} say \emph{lājṇe}
  \item \emph{Recall:} say the Marathi for to tidy up, then the Marathi for to notice, then say \emph{lājṇe} again
\end{itemize}

\begin{tcolorbox}[breakable,colback=teal!4,colframe=teal!35!black,title={Before you move on}]
What does \mr{लाजणे} mean? (\textquotedblleft{}To be shy, to blush\textquotedblright{}.) Say it once more.
\end{tcolorbox}
\section[\texorpdfstring{\mr{हरणे} (harṇe)}{हरणे (harṇe)}]{\mr{हरणे} (harṇe) — to be defeated, to lose (a game)}

\label{lesson:MR-C168-harne}



\begin{quote}
Before the new one: say the Marathi for to notice, then the Marathi for to be shy.
\end{quote}

\subsection*{You'll want to know: \mr{हरणे}}
\textbf{\mr{हरणे}} — \emph{harṇe} — \textquotedblleft{}to be defeated, to lose (a game)\textquotedblright{}.

\begin{grammarlens}[title={What you've built}]
One more word for this chapter.
\end{grammarlens}

\subsection*{Your turn}
\begin{itemize}
\raggedright
  \item \emph{Say it:} \emph{harṇe}
  \item \emph{Say it:} \emph{harṇe}, once more
  \item \emph{Say it:} say \emph{harṇe}
  \item \emph{Recall:} say the Marathi for to notice, then the Marathi for to be shy, then say \emph{harṇe} again
\end{itemize}

\begin{tcolorbox}[breakable,colback=teal!4,colframe=teal!35!black,title={Before you move on}]
What does \mr{हरणे} mean? (\textquotedblleft{}To be defeated, to lose (a game)\textquotedblright{}.) Say it once more.
\end{tcolorbox}
\section[\texorpdfstring{\mr{लपणे} (lapṇe)}{लपणे (lapṇe)}]{\mr{लपणे} (lapṇe) — to take cover, to hide}

\label{lesson:MR-C168-lapne}



\begin{quote}
Before the new one: say the Marathi for to be shy, then the Marathi for to be defeated.
\end{quote}

\subsection*{You'll want to know: \mr{लपणे}}
\textbf{\mr{लपणे}} — \emph{lapṇe} — \textquotedblleft{}to take cover, to hide\textquotedblright{}.

\begin{grammarlens}[title={What you've built}]
One more word for this chapter.
\end{grammarlens}

\subsection*{Your turn}
\begin{itemize}
\raggedright
  \item \emph{Say it:} \emph{lapṇe}
  \item \emph{Say it:} \emph{lapṇe}, once more
  \item \emph{Say it:} say \emph{lapṇe}
  \item \emph{Recall:} say the Marathi for to be shy, then the Marathi for to be defeated, then say \emph{lapṇe} again
\end{itemize}

\begin{tcolorbox}[breakable,colback=teal!4,colframe=teal!35!black,title={Before you move on}]
What does \mr{लपणे} mean? (\textquotedblleft{}To take cover, to hide\textquotedblright{}.) Say it once more.
\end{tcolorbox}
